\documentclass[10pt,a4paper]{amsart}
\usepackage{amsmath,amssymb,amsthm,graphicx}
% For commutative diagrams we use Xy-pic
\usepackage[all]{xy}

% We use 2cell for 2-commutative diagrams.
\xyoption{2cell}
% Theorem environments.
%
\theoremstyle{plain}
\newtheorem{theorem}[subsection]{Theorem}
\newtheorem{proposition}[subsection]{Proposition}
\newtheorem{lemma}[subsection]{Lemma}

\theoremstyle{definition}
\newtheorem{definition}[subsection]{Definition}
\newtheorem{example}[subsection]{Example}
\newtheorem{exercise}[subsection]{Exercise}
\newtheorem{situation}[subsection]{Situation}

\theoremstyle{remark}
\newtheorem{remark}[subsection]{Remark}
\newtheorem{remarks}[subsection]{Remarks}

\numberwithin{equation}{subsection}

% Macros
%
\def\lim{\mathop{\mathrm{lim}}\nolimits}
\def\colim{\mathop{\mathrm{colim}}\nolimits}
\def\Spec{\mathop{\mathrm{Spec}}}
\def\Hom{\mathop{\mathrm{Hom}}\nolimits}
\def\Ext{\mathop{\mathrm{Ext}}\nolimits}
\def\SheafHom{\mathop{\mathcal{H}\!\mathit{om}}\nolimits}
\def\SheafExt{\mathop{\mathcal{E}\!\mathit{xt}}\nolimits}
\def\Sch{\mathit{Sch}}
\def\Mor{\mathop{\mathrm{Mor}}\nolimits}
\def\Ob{\mathop{\mathrm{Ob}}\nolimits}
\def\Sh{\mathop{\mathit{Sh}}\nolimits}
\def\NL{\mathop{N\!L}\nolimits}
\def\CH{\mathop{\mathrm{CH}}\nolimits}
\def\proetale{{pro\text{-}\acute{e}tale}}
\def\etale{{\acute{e}tale}}
\def\QCoh{\mathit{QCoh}}
\def\Ker{\mathop{\mathrm{Ker}}}
\def\Im{\mathop{\mathrm{Im}}}
\def\Coker{\mathop{\mathrm{Coker}}}
\def\Coim{\mathop{\mathrm{Coim}}}

% Boxtimes
%
\DeclareMathSymbol{\boxtimes}{\mathbin}{AMSa}{"02}

%
% Macros for moduli stacks/spaces
%
\def\QCohstack{\mathcal{QC}\!\mathit{oh}}
\def\Cohstack{\mathcal{C}\!\mathit{oh}}
\def\Spacesstack{\mathcal{S}\!\mathit{paces}}
\def\Quotfunctor{\mathrm{Quot}}
\def\Hilbfunctor{\mathrm{Hilb}}
\def\Curvesstack{\mathcal{C}\!\mathit{urves}}
\def\Polarizedstack{\mathcal{P}\!\mathit{olarized}}
\def\Complexesstack{\mathcal{C}\!\mathit{omplexes}}
% \Pic is the operator that assigns to X its picard group, usage \Pic(X)
% \Picardstack_{X/B} denotes the Picard stack of X over B
% \Picardfunctor_{X/B} denotes the Picard functor of X over B
\def\Pic{\mathop{\mathrm{Pic}}\nolimits}
\def\Picardstack{\mathcal{P}\!\mathit{ic}}
\def\Picardfunctor{\mathrm{Pic}}
\def\Deformationcategory{\mathcal{D}\!\mathit{ef}}
\usepackage{cleveref}
\crefname{originallemma}{original lemma}{original lemmas}
\Crefname{originallemma}{Original lemma}{Original lemmas}
\makeatletter
\newlabel{sites-cohomology-lemma-kill-cohomology-class-on-covering}{{01FW}{source}{}{sourcepointer-1}{}}
\newlabel{sites-cohomology-lemma-kill-cohomology-class-on-covering@cref}{{[originallemma][0][]01FW}{[1][0][]source}{}{}{}}
\makeatother
\title{Locality of cohomology: a cohomology class on a site dies on a covering}
\author{The Stacks Project Authors}
\date{stacks-project:etale-cohomology.tex:local-snapshot-20261011; GNU FDL 1.2 or later}
\begin{document}
\maketitle
\noindent\textit{Editorial note.} All delivered mathematical text is the author's own, copied from the pinned Stacks Project chapter file etale-cohomology.tex; no statement or proof step was rewritten, reordered or removed. The author writes 'is exact in' where 'is exact at' is meant; this is reproduced verbatim and not corrected.\par\medskip
\noindent\textit{Editorial note.} The preamble is the chapter's own native preamble excerpt used for this source: the Xy-pic setup at lines 14-18 and the theorem declarations with the mathematical macros at lines 154-217 of sources/stacks-etale-cohomology/preamble.tex. The quoted class line and the xr-hyper externaldocument block are not included because the assembler supplies a plain amsart wrapper; no macro was invented or redefined.\par\medskip
\noindent\textit{Editorial note.} The cross-chapter pointer is located in the project's own tag list staged in the snapshot (sources/stacks-chow/tags.txt), at the verbatim line carrying its canonical tag.\par\medskip
\noindent\textit{Allowed standard prerequisites (editorial).}\begin{itemize}
\item Cohomology of an abelian sheaf on a site computed from an injective resolution. Standard in this chapter: H\textasciicircum{}p(U, F) is the p-th cohomology of the complex of sections of an injective resolution F -> I\textasciicircum{}bullet. Supplies the representative of the class and the fact that restriction is computed on representatives; the covering argument is the counted content.
\item Exactness of the resolution in positive degrees means local triviality of sections of the kernel. Standard sheaf theory: I\textasciicircum{}\{p-1\} -> I\textasciicircum{}p -> I\textasciicircum{}\{p+1\} exact in Ab(C) means a section over U killed by d becomes a coboundary after restriction to a covering. Supplies the existence of the local primitives xi\_i; identifying them and concluding the vanishing is inside the counted proof.
\end{itemize}\medskip
\ifdefined\hypertarget\hypertarget{sourcepointer-1}{}\fi
\noindent\textit{Source reference.} Original lemma 01FW: sources/stacks-chow/tags.txt, source lines 1885--1885. Cross-chapter 'for more details' pointer at proof line 2276, invoked after the argument is complete; the resolution, the local primitives and the vanishing are all retained verbatim in the counted span, so this pointer replaces nothing.\par\medskip
\begin{lemma}[Locality of cohomology]
\label{lemma-locality-cohomology}
Let $\mathcal{C}$ be a site, $\mathcal{F}$ an abelian sheaf on $\mathcal{C}$,
$U$ an object of $\mathcal{C}$, $p > 0$ an integer and $\xi \in
H^p(U, \mathcal{F})$. Then there exists a covering
$\mathcal{U} = \{U_i \to U\}_{i \in I}$ of $U$ in $\mathcal{C}$
such that $\xi |_{U_i} = 0$ for all $i \in I$.
\end{lemma}

\begin{proof}
Choose an injective resolution $\mathcal{F} \to \mathcal{I}^\bullet$. Then
$\xi$ is represented by a cocycle $\tilde{\xi} \in \mathcal{I}^p(U)$
with $d^p(\tilde{\xi}) = 0$. By assumption, the sequence
$\mathcal{I}^{p - 1} \to \mathcal{I}^p \to \mathcal{I}^{p + 1}$ in exact in
$\textit{Ab}(\mathcal{C})$, which means that there exists a covering
$\mathcal{U} = \{U_i \to U\}_{i \in I}$ such that
$\tilde{\xi}|_{U_i} = d^{p - 1}(\xi_i)$ for some
$\xi_i \in \mathcal{I}^{p-1}(U_i)$. Since
the cohomology class $\xi|_{U_i}$ is represented by the cocycle
$\tilde{\xi}|_{U_i}$ which is a coboundary, it vanishes.
For more details see
Cohomology on Sites,
Lemma \ref{sites-cohomology-lemma-kill-cohomology-class-on-covering}.
\end{proof}

\end{document}
